\documentclass{article}
\usepackage{booktabs}
\usepackage{tabularx}
\usepackage[margin=1in]{geometry}
\begin{document}

\begin{table}[tbp] \centering
\newcolumntype{R}{>{\raggedleft\arraybackslash}X}
\newcolumntype{L}{>{\raggedright\arraybackslash}X}
\newcolumntype{C}{>{\centering\arraybackslash}X}

\caption{__000006}
\label{tab:cleavage_index}
\begin{tabularx}{\linewidth}{@{}lCCCCC@{}}

\toprule
&{(1)}&{(2)}&{(3)}&{(4)}&{(5)} \tabularnewline \midrule
{Vote no.}&{Year}&{Title (short)}&{Lang gap (pp)}&{Rho (between-lang)}&{} \tabularnewline
\midrule \addlinespace[\belowrulesep]
56&1900&Gesetz zur Kranken-, Unfall- und Militärversicher&--15.8&0.282&
57&1900&Initiative «für die Proporzwahl des Nationalrate&14.2&0.096&
58&1900&Initiative für die Volkswahl des Bundesrates&--0.5&0.000&
59&1902&Unterstützung der Primarschule durch den Bund&12.9&0.152&
60&1903&Zolltarifgesetz&--14.7&0.080&
61&1903&Bundesstrafrecht (Anstiftung Militärpflichtiger z&3.1&0.011&
62&1903&Initiative «für die Wahl des Nationalrates aufgr&6.9&0.020&
63&1903&Artikel über die Regulierung des Alkoholhandels&7.3&0.073&\$\\dagger\$
64&1905&Ausdehnung des Erfinderschutzes&14.7&0.176&
65&1906&Lebensmittelgesetz&--2.7&0.003&\$\\dagger\$
66&1907&Militärorganisation der schweizerischen Eidgenoss&--12.2&0.139&
67&1908&Verfassungsartikel zur Gesetzgebung über das Gewe&--4.4&0.026&
68&1908&Initiative «für ein Absinthverbot»&--17.1&0.343&\$\\dagger\$
69&1908&Verfassungsartikel über die Wasserkraft und Elekt&8.7&0.079&
70&1910&Initiative «für die Proporzwahl des Nationalrate&--3.0&0.006&
\bottomrule \addlinespace[\belowrulesep]

\end{tabularx}
\\ \parbox{\linewidth}{\footnotesize Notes: Language Cleavage Index across federal referenda, 1900-1910. The Language Gap (column 4) is the difference in mean yes-vote share between French- and German-majority cantons (defined by french\\_share >= 0.5; threshold definition with 5 French / 20 German cantons; this is DISTINCT from the strict \{VD, VS, NE, GE\} canton-set definition used in the turnout-mobilization analysis Table 20). The Between-Language Variance Share rho (column 5) is the share of total cross-cantonal variance in yes\\_pct attributable to the language partition (rho = (n\\_fr (y\\_fr - y\\_grand)\^{}2 + n\\_ge (y\\_ge - y\\_grand)\^{}2) / (25 var\\_total\\_yes)). Wine-/alcohol-relevant votes (\\#63 federal alcohol-trade regulation 1903 [null reference]; \\#65 Lebensmittelgesetz 1906; \\#68 Absinthverbot 1908) are marked with daggers. Per-vote rho for the three wine-/alcohol-relevant votes: rho\\_63 = 0.073, rho\\_65 = 0.003, rho\\_68 = 0.343. Mean rho across wine-/alcohol-relevant votes (n=3) = 0.139 vs mean rho across other 12 votes = 0.089 (two-sample ttest p = 0.472 under the threshold definition). Robustness with NE + GE excluded (N=23 cantons; the two French cantons that rejected vote \\#68): mean rho | wine\\_relevant = 0.182 vs mean rho | other = 0.097. Substantive heterogeneity: the cleavage-attenuation hypothesis is CONFIRMED for vote \\#65 (rho\\_65 very low, the food law attracted broad cross-language coalition support; consistent with wine-industry interests pulling across the cleavage on a vote about wine adulteration / substitute beverages) but is NOT confirmed for vote \\#68 (rho\\_68 elevated, the absinthe ban produced a sharp language-aligned vote pattern). The likely structural reason: under the threshold definition the French-canton set IS the wine-canton set (VD, VS, NE, GE produce nearly all Swiss wine), so on the absinthe vote where wine cantons vote yes, the language partition coincides with the wine partition rather than cutting across it. The strategist's prediction was that German wine cantons (SH, ZH, AG, TG) would vote with French wine cantons and pull the German-canton mean upward, attenuating rho; empirically, those German wine cantons are too few or too small relative to the German-canton denominator (20 cantons) to produce that attenuation on \\#68. Vote \\#65 shows the predicted attenuation because the food law's coalition was broader than the wine industry alone (public-health support cross-cut the language line). Framing: the cleavage index documents language as the dominant cultural-political cleavage in 1900-1910 Swiss federal voting (consistent with the Gelbach decomposition in Table 15, LANG channel = 99 percent of the headline Simpson sign-flip), and the heterogeneity in rho across the three alcohol-/wine-relevant votes provides texture on coalition structure: \\#65 had a cross-cutting coalition; \\#68 had a language-aligned coalition. We do NOT interpret either pattern as direct evidence of the moralist-temperance coalition mechanism (formal H3 test in Table 17 was null at N=25 under both spec variants); the cleavage-index heterogeneity is consistent with multiple mechanisms and the data at this sample size cannot uniquely identify which mechanism drives each vote's coalition structure.}
\end{table}
\end{document}
